% MT066 / explanation — canonical editable content.
损失\(\mathcal L\)衡量预测错得多严重。偏导数\(\partial\mathcal L/\partial w\)表示只把权重\(w\)调大一点时，损失会怎样变；所有参数的偏导数组成梯度。训练按\(w\leftarrow w-\eta\,\partial\mathcal L/\partial w\)更新，\(\eta>0\)是控制每步幅度的学习率。传到前面某层的梯度若很小，该层每次更新也很小，训练就很慢，这叫梯度消失，不要求梯度严格等于0。


\textbf{为什么深度会放大这个问题？} 链式法则把一条路径上的局部变化率相乘。设共有\(L\)层，\(h_l\)表示第\(l\)层输出，\(w_l,b_l\)分别是该层权重与偏置。对于一条只有标量相连的计算链，\(h_l=\sigma(w_lh_{l-1}+b_l)\)，有



\[
\frac{\partial h_L}{\partial h_0}
=\prod_{l=1}^{L}\left[w_l\,\sigma'(w_lh_{l-1}+b_l)\right].
\]


每一项都在回答“前一层动一点，后一层动多少”。若每层缩成原来的0.2，四层后剩\(0.2^4=0.0016\)。路径越长，多个小于1的因子越可能把梯度压得很小；矩阵网络还需要把多条路径贡献相加，但仍有连乘机制。


\textbf{Sigmoid为何容易带来小因子？} \(\sigma(t)=1/(1+e^{-t})\)把输入压到0与1之间，\(e\approx2.718\)是自然指数的底数。其导数为\(\sigma(t)[1-\sigma(t)]\)。令\(p=\sigma(t)\)，则\(p(1-p)=1/4-(p-1/2)^2\le1/4\)，两端饱和时还更接近0。以单位权重为例，经过\(L\)个sigmoid的单路径梯度不超过\((1/4)^L\)。一般网络权重也参与乘积，因此不能只数激活层数就断言必然消失。


\textbf{ReLU缓解了哪一段？} ReLU为\(\max(0,t)\)，在\(t>0\)时斜率1，活跃路径不会再被sigmoid的小斜率压缩；在\(t<0\)时斜率0，仍可能让某些单元没有梯度。权重过小或其他层的因素也仍存在，所以ReLU是缓解手段，不是万能保证。


\textbf{出处：} 2025A期中 Q10，原卷第7页（10分）


\textbf{答案依据：} 2025A期中 Q10 印刷题干；本题解为我们的推导，不采用扫描件的学生作答。
